\documentclass[UTF8]{ctexart}
\usepackage{geometry}
\usepackage{hyperref}
\usepackage{listings}
\usepackage{xcolor}

\geometry{a4paper, left=2.5cm, right=2.5cm, top=2.5cm, bottom=2.5cm}

\title{Agent开发周报 260318}
\author{刘德辰}
\date{2026年3月18日}

\begin{document}

\maketitle

\section{车辆报告与线路报告收口：三类报告能力对齐}

\textbf{背景}：此前驾驶员报告已完成主链路稳定化，但车辆报告和线路报告在无数据返回时的处理方式与驾驶员报告不一致，报告生成后的追问也存在被路由层自行截答、未能继续走业务咨询路径的问题。三类报告的结构化输出格式也各自为政，维护成本高。

\textbf{方案说明}：以驾驶员报告的成熟路径为基准，将车辆报告和线路报告的无数据兜底、追问路由和输出格式统一收口，确保三类报告在同等场景下行为一致。

\textbf{功能说明}：
\begin{itemize}
\item 车辆报告补齐无数据时的兜底回复，不再返回空内容或异常，与驾驶员报告行为对齐
\item 线路报告经三轮迭代，完成字段完整性、关键指标映射和输出格式收敛
\item 报告生成后的追问（如"为什么这个风险排第一""给我几条管理动作"）统一改走业务咨询路径，不再由路由层自行回答
\item 三类报告的 JSON 输出解析和格式校验逻辑统一维护，减少各类报告之间的行为差异
\end{itemize}

\section{规则编辑体验优化：操作意图识别与冲突兜底}

\textbf{背景}：用户在编辑已有规则时，常见的"补充关键点""再加一条"等表达，系统此前会将原有数组内容整体替换，而不是在已有内容上追加；遇到模糊诉求（如"更适合早班司机"）也会强行映射到字段值，导致规则被改乱；同一轮出现矛盾要求时，系统也会将冲突内容一并写入，产生前后矛盾的规则。

\textbf{方案说明}：在对话式规则收集与规则编译两个环节同步增强意图识别能力，区分"追加"与"替换"，并在出现模糊诉求或冲突指令时主动澄清而非强行执行。

\textbf{功能说明}：
\begin{itemize}
\item 编辑态下，"补上""增加""追加""再加一个"等表达默认理解为在原有内容上追加，不会清空已有字段
\item 收到抽象诉求（如"再严谨些""口径统一一点"）时，系统不再擅自猜测字段映射，而是提出单一澄清问题
\item 同一轮出现相互矛盾的要求时，系统不会将冲突内容一并写入，会先暂停并引导用户确认意图
\item 规则编译环节同步感知编辑上下文，能正确区分本轮是追加还是整体替换
\end{itemize}

\section{规则配置服务独立化：与主链路解耦}

\textbf{背景}：规则配置涉及多轮对话、草稿管理、状态流转和最终编译多个环节，此前所有逻辑都集中在系统运行时主文件中，每次迭代规则配置功能都需要在一个几千行的文件里定位和修改，维护成本高、出错风险大。

\textbf{方案说明}：将规则配置相关的状态机、草稿存储、业务逻辑和 HTTP 接口全部迁移到独立的 \texttt{rule-config} 领域模块，与主链路完全解耦，后续迭代可以独立进行。

\textbf{功能说明}：
\begin{itemize}
\item 规则配置的类型定义、纯逻辑函数、会话状态机、草稿存储和服务主体各自独立成文件
\item HTTP 接口层与对话入口层分离，互不干扰
\item 从主运行时文件中删除超过 1100 行的规则配置旧代码，运行时主文件职责进一步收窄
\item 拆分过程分七次提交递进完成，每次提交保持编译与功能可用，无回归
\end{itemize}

\section{报表查找解析统一化：实体识别与候选澄清}

\textbf{背景}：用户在查询驾驶员、车辆、线路、事故报告时，经常出现同名、编号模糊、输入不完整等情况，此前各类报告的实体查找逻辑分散，处理策略不统一，部分场景下会直接报错或生成错误报告。

\textbf{方案说明}：将四类报告实体的查找与解析逻辑统一收口，明确区分"精确匹配"与"需要澄清"两种结果，澄清场景进一步细化原因，供系统生成针对性的引导回复。

\textbf{功能说明}：
\begin{itemize}
\item 驾驶员、车辆、线路、事故四类实体查找逻辑统一管理，行为一致
\item 查找结果明确分为"已确定"和"需澄清"两类，需澄清时标明原因（输入过短、多个候选、同名、未找到等）
\item 路由层依据查找结果决定是直接生成报告还是先向用户确认，减少因实体误判生成错误报告的情况
\end{itemize}

\section{MCP 外部工具接入：网络链路打通与工具质量评估}

\textbf{背景}：交信投的车辆管理 MCP 服务部署在内网（\texttt{10.181.92.x}），需要通过 VPN 才能访问。此前链路未打通，工具描述质量也未经评估，无法实际接入。

\textbf{方案说明}：分两步推进——先打通从公网到内网的完整访问链路，再对 MCP 服务暴露的工具进行可达性验证和描述质量审计。

\textbf{功能说明}：
\begin{itemize}
\item 在 VPS 上配置 OpenVPN 客户端为 systemd 持久化服务，打通 VPS 到内网 \texttt{10.181.92.0/24} 的路由，链路验证通过（HTTP 405 确认服务可达）
\item 通过 Cloudflare Tunnel 将内网 MCP 服务映射到公网域名 \texttt{buso-rpformcp.canocache.com}，完成 CF Workers → cloudflared → VPS → OpenVPN → 内网 MCP 的完整链路
\item 编写 \texttt{mcp\_probe.py}，验证 MCP 初始化握手（\texttt{initialize} → \texttt{notifications/initialized} → \texttt{tools/list}）全流程可达，协议版本 \texttt{2025-03-26} 握手成功，获取到 8 个车辆管理工具
\item 编写 \texttt{mcp\_tool\_audit.py}，对 8 个工具逐一进行描述清晰度评估，发现所有工具存在描述重复工具名、参数说明仅翻译字段名、写操作缺少 \texttt{required} 声明三类系统性问题
\item 输出 \texttt{mcp-tool-description-guide.md} 规范文档，说明工具描述对 Agent 决策的影响机制及五条具体修复规则，已发送交信投团队跟进改进
\end{itemize}

\section{Agent 主链路架构持续瘦身}

\textbf{背景}：Agent 主运行时文件随功能迭代不断膨胀，目前仍承载着会话编排、路由分发、认证管理等多项职责，给日常维护和并行开发带来较大摩擦。

\textbf{方案说明}：延续上周的分层拆分方向，本周重点将会话处理、路由分发和认证状态管理迁出，Runtime 进一步退为薄入口层。

\textbf{功能说明}：
\begin{itemize}
\item 会话请求处理主链路（规则路由、Worker 执行、流式回复、消息存储）提取为独立服务
\item 路由分发逻辑（Omni 决策、场景匹配、报告追问检测）提取为独立服务
\item 认证上下文构建和会话 Cookie 管理提取为基础设施层独立模块
\item LLM 调用层封装为独立模块，统一类型定义，减少各业务模块对底层调用细节的直接依赖
\end{itemize}

\section{技术债务与限制}

\begin{itemize}
\item \textbf{报告追问路由仍有漏案例}：已发现在部分多轮报告对话中，路由层仍会直接回答业务问题而非转交业务咨询路径，根因为路由判断对上下文感知不足，已记录待修复
\item \textbf{主运行时拆分仍在进行}：研究评估、情景、会话等领域的逻辑尚未迁出，后续继续推进
\item \textbf{MCP 接入仍处于基础阶段}：当前参数描述覆写为系统侧临时方案，外部服务团队的工具描述质量改善需持续跟进
\item \textbf{规则模板库}：建立常用规则模板，降低配置成本（无限期延后）。
\end{itemize}

\section{下周计划}

\begin{enumerate}
\item 报告追问路由修复：针对已记录的路由漏案例，约束路由判断在报告上下文下的行为，确保建议类问题稳定走业务咨询路径
\item 主链路拆分第三阶段：继续迁出研究评估、情景匹配等领域逻辑
\item MCP 接入联调：周五上班前接入MCP tools，设定工具调用策略，在真实工具上跑部分测试。
\item 三类报告回归覆盖：补充驾驶员、车辆、线路报告的回归用例，覆盖正常生成、无数据兜底和追问路由三类场景
\end{enumerate}

\end{document}